% Options for packages loaded elsewhere
\PassOptionsToPackage{unicode}{hyperref}
\PassOptionsToPackage{hyphens}{url}
\documentclass[
]{article}
\usepackage{xcolor}
\usepackage[margin=1in]{geometry}
\usepackage{amsmath,amssymb}
\setcounter{secnumdepth}{-\maxdimen} % remove section numbering
\usepackage{iftex}
\ifPDFTeX
  \usepackage[T1]{fontenc}
  \usepackage[utf8]{inputenc}
  \usepackage{textcomp} % provide euro and other symbols
\else % if luatex or xetex
  \usepackage{unicode-math} % this also loads fontspec
  \defaultfontfeatures{Scale=MatchLowercase}
  \defaultfontfeatures[\rmfamily]{Ligatures=TeX,Scale=1}
\fi
\usepackage{lmodern}
\ifPDFTeX\else
  % xetex/luatex font selection
\fi
% Use upquote if available, for straight quotes in verbatim environments
\IfFileExists{upquote.sty}{\usepackage{upquote}}{}
\IfFileExists{microtype.sty}{% use microtype if available
  \usepackage[]{microtype}
  \UseMicrotypeSet[protrusion]{basicmath} % disable protrusion for tt fonts
}{}
\makeatletter
\@ifundefined{KOMAClassName}{% if non-KOMA class
  \IfFileExists{parskip.sty}{%
    \usepackage{parskip}
  }{% else
    \setlength{\parindent}{0pt}
    \setlength{\parskip}{6pt plus 2pt minus 1pt}}
}{% if KOMA class
  \KOMAoptions{parskip=half}}
\makeatother
\usepackage{color}
\usepackage{fancyvrb}
\newcommand{\VerbBar}{|}
\newcommand{\VERB}{\Verb[commandchars=\\\{\}]}
\DefineVerbatimEnvironment{Highlighting}{Verbatim}{commandchars=\\\{\}}
% Add ',fontsize=\small' for more characters per line
\usepackage{framed}
\definecolor{shadecolor}{RGB}{248,248,248}
\newenvironment{Shaded}{\begin{snugshade}}{\end{snugshade}}
\newcommand{\AlertTok}[1]{\textcolor[rgb]{0.94,0.16,0.16}{#1}}
\newcommand{\AnnotationTok}[1]{\textcolor[rgb]{0.56,0.35,0.01}{\textbf{\textit{#1}}}}
\newcommand{\AttributeTok}[1]{\textcolor[rgb]{0.13,0.29,0.53}{#1}}
\newcommand{\BaseNTok}[1]{\textcolor[rgb]{0.00,0.00,0.81}{#1}}
\newcommand{\BuiltInTok}[1]{#1}
\newcommand{\CharTok}[1]{\textcolor[rgb]{0.31,0.60,0.02}{#1}}
\newcommand{\CommentTok}[1]{\textcolor[rgb]{0.56,0.35,0.01}{\textit{#1}}}
\newcommand{\CommentVarTok}[1]{\textcolor[rgb]{0.56,0.35,0.01}{\textbf{\textit{#1}}}}
\newcommand{\ConstantTok}[1]{\textcolor[rgb]{0.56,0.35,0.01}{#1}}
\newcommand{\ControlFlowTok}[1]{\textcolor[rgb]{0.13,0.29,0.53}{\textbf{#1}}}
\newcommand{\DataTypeTok}[1]{\textcolor[rgb]{0.13,0.29,0.53}{#1}}
\newcommand{\DecValTok}[1]{\textcolor[rgb]{0.00,0.00,0.81}{#1}}
\newcommand{\DocumentationTok}[1]{\textcolor[rgb]{0.56,0.35,0.01}{\textbf{\textit{#1}}}}
\newcommand{\ErrorTok}[1]{\textcolor[rgb]{0.64,0.00,0.00}{\textbf{#1}}}
\newcommand{\ExtensionTok}[1]{#1}
\newcommand{\FloatTok}[1]{\textcolor[rgb]{0.00,0.00,0.81}{#1}}
\newcommand{\FunctionTok}[1]{\textcolor[rgb]{0.13,0.29,0.53}{\textbf{#1}}}
\newcommand{\ImportTok}[1]{#1}
\newcommand{\InformationTok}[1]{\textcolor[rgb]{0.56,0.35,0.01}{\textbf{\textit{#1}}}}
\newcommand{\KeywordTok}[1]{\textcolor[rgb]{0.13,0.29,0.53}{\textbf{#1}}}
\newcommand{\NormalTok}[1]{#1}
\newcommand{\OperatorTok}[1]{\textcolor[rgb]{0.81,0.36,0.00}{\textbf{#1}}}
\newcommand{\OtherTok}[1]{\textcolor[rgb]{0.56,0.35,0.01}{#1}}
\newcommand{\PreprocessorTok}[1]{\textcolor[rgb]{0.56,0.35,0.01}{\textit{#1}}}
\newcommand{\RegionMarkerTok}[1]{#1}
\newcommand{\SpecialCharTok}[1]{\textcolor[rgb]{0.81,0.36,0.00}{\textbf{#1}}}
\newcommand{\SpecialStringTok}[1]{\textcolor[rgb]{0.31,0.60,0.02}{#1}}
\newcommand{\StringTok}[1]{\textcolor[rgb]{0.31,0.60,0.02}{#1}}
\newcommand{\VariableTok}[1]{\textcolor[rgb]{0.00,0.00,0.00}{#1}}
\newcommand{\VerbatimStringTok}[1]{\textcolor[rgb]{0.31,0.60,0.02}{#1}}
\newcommand{\WarningTok}[1]{\textcolor[rgb]{0.56,0.35,0.01}{\textbf{\textit{#1}}}}
\usepackage{graphicx}
\makeatletter
\newsavebox\pandoc@box
\newcommand*\pandocbounded[1]{% scales image to fit in text height/width
  \sbox\pandoc@box{#1}%
  \Gscale@div\@tempa{\textheight}{\dimexpr\ht\pandoc@box+\dp\pandoc@box\relax}%
  \Gscale@div\@tempb{\linewidth}{\wd\pandoc@box}%
  \ifdim\@tempb\p@<\@tempa\p@\let\@tempa\@tempb\fi% select the smaller of both
  \ifdim\@tempa\p@<\p@\scalebox{\@tempa}{\usebox\pandoc@box}%
  \else\usebox{\pandoc@box}%
  \fi%
}
% Set default figure placement to htbp
\def\fps@figure{htbp}
\makeatother
\setlength{\emergencystretch}{3em} % prevent overfull lines
\providecommand{\tightlist}{%
  \setlength{\itemsep}{0pt}\setlength{\parskip}{0pt}}
\usepackage{bookmark}
\IfFileExists{xurl.sty}{\usepackage{xurl}}{} % add URL line breaks if available
\urlstyle{same}
\hypersetup{
  pdftitle={Hamiltonian Monte Carlo},
  hidelinks,
  pdfcreator={LaTeX via pandoc}}

\title{Hamiltonian Monte Carlo}
\author{}
\date{\vspace{-2.5em}2026-04-17}

\begin{document}
\maketitle

\subsection{Introduction}\label{introduction}

Hamiltonian Monte Carlo (HMC) is the current state of the art for
sampling from continuous Bayesian posteriors of moderate-to-high
dimension. Where random-walk Metropolis takes small, undirected steps
and Gibbs cycles through coordinates, HMC uses gradient information to
propose long-distance moves that almost always land in regions of high
posterior density. The No-U-Turn Sampler (NUTS), Stan's default,
automates the choice of trajectory length so that users do not need to
tune step counts by hand. The result is mixing rates that are orders of
magnitude better than coordinate-wise samplers for posteriors with
strong correlations.

\subsection{Prerequisites}\label{prerequisites}

A working understanding of MCMC, basic calculus (gradients), and the
relationship between log-posterior gradients and the geometry of the
posterior surface. Familiarity with Stan or \texttt{brms} makes the
implementation step trivial.

\subsection{Theory}\label{theory}

HMC introduces auxiliary momentum variables \(p\) and samples from the
joint distribution
\(p(\theta, p) = p(\theta \mid y) \mathcal N(p \mid 0, M)\), where \(M\)
is the mass matrix. Hamiltonian dynamics conserve the joint Hamiltonian
\(H(\theta, p) = -\log p(\theta \mid y) + \tfrac{1}{2} p^T M^{-1} p\),
so trajectories simulated with leapfrog integration drift through level
sets of constant log-density and explore the posterior efficiently.
After \(L\) leapfrog steps with step size \(\epsilon\), the proposal is
accepted with a Metropolis correction that compensates for integrator
error. NUTS replaces the manual choice of \(L\) with a ``no-U-turn''
criterion that terminates the trajectory when it starts doubling back,
eliminating the most error-prone tuning parameter of plain HMC.

\subsection{Assumptions}\label{assumptions}

A continuous, almost-everywhere differentiable log-posterior. Discrete
parameters must be marginalised out. The mass matrix should approximate
the posterior covariance --- Stan adapts it during warm-up --- and the
step size must be small enough to keep integrator error in check.

\subsection{R Implementation}\label{r-implementation}

\begin{Shaded}
\begin{Highlighting}[]
\FunctionTok{library}\NormalTok{(brms)}

\FunctionTok{set.seed}\NormalTok{(}\DecValTok{2026}\NormalTok{)}
\NormalTok{d }\OtherTok{\textless{}{-}} \FunctionTok{data.frame}\NormalTok{(}\AttributeTok{x =} \FunctionTok{rnorm}\NormalTok{(}\DecValTok{100}\NormalTok{))}
\NormalTok{d}\SpecialCharTok{$}\NormalTok{y }\OtherTok{\textless{}{-}} \DecValTok{2} \SpecialCharTok{+} \FloatTok{1.5} \SpecialCharTok{*}\NormalTok{ d}\SpecialCharTok{$}\NormalTok{x }\SpecialCharTok{+} \FunctionTok{rnorm}\NormalTok{(}\DecValTok{100}\NormalTok{)}

\NormalTok{fit }\OtherTok{\textless{}{-}} \FunctionTok{brm}\NormalTok{(y }\SpecialCharTok{\textasciitilde{}}\NormalTok{ x, }\AttributeTok{data =}\NormalTok{ d,}
           \AttributeTok{chains =} \DecValTok{4}\NormalTok{, }\AttributeTok{iter =} \DecValTok{2000}\NormalTok{, }\AttributeTok{refresh =} \DecValTok{0}\NormalTok{,}
           \AttributeTok{control =} \FunctionTok{list}\NormalTok{(}\AttributeTok{adapt\_delta =} \FloatTok{0.95}\NormalTok{, }\AttributeTok{max\_treedepth =} \DecValTok{10}\NormalTok{))}

\FunctionTok{summary}\NormalTok{(fit)}
\CommentTok{\# Diagnostics: R{-}hat, divergences, energy}
\NormalTok{rstan}\SpecialCharTok{::}\FunctionTok{check\_hmc\_diagnostics}\NormalTok{(fit}\SpecialCharTok{$}\NormalTok{fit)}
\end{Highlighting}
\end{Shaded}

\subsection{Output \& Results}\label{output-results}

\texttt{summary(fit)} reports posterior summaries with \(\hat R\), ESS,
and credible intervals. \texttt{rstan::check\_hmc\_diagnostics()}
reports divergent transitions, tree-depth saturation, energy diagnostics
(E-BFMI), and is the first-line check that the geometry was navigable.
Healthy HMC delivers \(\hat R\) near 1.00, no divergent transitions, and
ESS that is a substantial fraction of the post-warm-up draws.

\subsection{Interpretation}\label{interpretation}

A reporting sentence: ``HMC/NUTS produced 4,000 post-warm-up draws
across four chains with no divergent transitions and \(\hat R = 1.00\)
for every parameter; the slope posterior was 1.52 (95 \% credible
interval 1.31 to 1.74).'' When divergences occur, report them and
describe the remedy (reparameterisation, raising \texttt{adapt\_delta});
a divergent transition is a missed region of the posterior, not a tuning
curiosity.

\subsection{Practical Tips}\label{practical-tips}

\begin{itemize}
\tightlist
\item
  For hierarchical models with funnel geometry, switch to non-centred
  parameterisation; \texttt{brms} does this automatically when you
  specify \texttt{prior(...\ ,\ class\ =\ "sd",\ ...)}.
\item
  Raising \texttt{adapt\_delta} to 0.95 or 0.99 reduces divergences at
  the cost of smaller step sizes and slower sampling; reparameterisation
  is usually a better fix.
\item
  Tree-depth saturation (\texttt{treedepth\_\_} hitting
  \texttt{max\_treedepth}) signals that NUTS could not turn around
  within the budget --- increase \texttt{max\_treedepth} or simplify the
  model.
\item
  Reparameterise constrained variables onto unconstrained scales when
  you can: \texttt{log(sigma)} rather than
  \texttt{sigma\ \textgreater{}\ 0} keeps the leapfrog integrator
  stable.
\item
  HMC scales to thousands of parameters where Gibbs or random-walk MH
  would be unworkable, but it cannot rescue an ill-conditioned
  posterior; geometry matters more than algorithm.
\end{itemize}

\subsection{R Packages Used}\label{r-packages-used}

\texttt{rstan} and \texttt{cmdstanr} as the HMC/NUTS backends,
\texttt{brms} and \texttt{rstanarm} for high-level interfaces,
\texttt{bayesplot} for diagnostic plots (\texttt{mcmc\_pairs},
\texttt{mcmc\_parcoord}) that visualise divergent transitions, and
\texttt{posterior} for the underlying draw arrays.

\subsection{For Reviewers}\label{for-reviewers}

\textbf{What to look for in a paper using this method.}

\begin{itemize}
\tightlist
\item
  \textbf{Common misapplications.}
\item
  \textbf{Diagnostics that should be reported but often aren't.}
\item
  \textbf{Red flags in tables and figures.}
\item
  \textbf{What to verify.}
\item
  \textbf{What an adequate Methods paragraph must contain.}
\end{itemize}

\subsection{See also --- labs in this
chapter}\label{see-also-labs-in-this-chapter}

\begin{itemize}
\tightlist
\item
  \href{labs/lab-bayesian-thinking-control-vs-decision-error.Rmd}{Bayesian
  thinking; α control vs decision error}
\item
  \href{labs/lab-brms-stan-loo-hierarchical-models.Rmd}{brms / Stan,
  LOO, hierarchical models}
\end{itemize}

\end{document}
